\section{Cómo leer el resto de CS25: cuatro interfaces y una lista de verificación}

El valor final de la introducción no es detener la historia en GPT-3/BERT, sino ofrecer un lenguaje de comparación para el resto del curso. Ya sea que el ponente hable de visión, RL, modelos generativos, biology o modelos de lenguaje más grandes, todo puede desglosarse con cuatro preguntas: token interface, interaction pattern, training objective y deployment contract.

\subsection{Método de lectura de las cuatro interfaces}

\begin{importantbox}{Cada sesión responde cuatro preguntas}
\begin{enumerate}
  \item \textbf{Token interface}: ¿en qué unidades se divide la entrada y cómo se representan la position y la modality?
  \item \textbf{Interaction pattern}: ¿full, local, causal, cross, sparse o recurrent?
  \item \textbf{Training objective}: ¿next-token, mask, contrastive, control, ranking o multi-task?
  \item \textbf{Deployment contract}: ¿la salida se usa para generation, representation, decision, retrieval o action? ¿Cómo se evalúan las fallas?
\end{enumerate}
\end{importantbox}

\teachervoice{Al final, los tres instructores dirigen al público hacia la serie completa de guest-speakers. Nota de clase: la introducción solo aporta un vocabulary común; el aprendizaje real ocurre al ver cómo los ponentes siguientes modifican el token, la mask, el objective y el system boundary.}

\begin{knowledgebox}{Plantilla de pregunta de investigación}
Primero elige una falla del dominio: contexto largo, pocos datos, latencia en tiempo real, decisiones causales o restricciones de seguridad; después señala qué interfaz del Transformer estándar causa la limitación; propone un cambio verificable; presenta el baseline, la ablation y el presupuesto de recursos; por último, explica qué no pueden demostrar las figuras y tablas. Así, las «aplicaciones de Transformer» dejan de ser un simple cambio de conjunto de datos y se convierten en diseño de investigación.
\end{knowledgebox}

\subsection{Resumen de la sección}

El resto de CS25 no debe leerse como un catálogo de modelos. El método de las cuatro interfaces ayuda a distinguir el esqueleto general, la adaptación al dominio, la señal de entrenamiento y la responsabilidad del producto, y convierte cada guest talk en un diseño de sistema comparable.
